\documentclass[10pt]{beamer}
\usepackage{tikz}
\usetikzlibrary{positioning}

\usepackage[diamond]{beamertheme-geode}

\title{Crystal Structures, Briefly}
\author{Dana Whitfield}
\date{\today}

\begin{document}

\geotitle

\section{Why Crystals Matter}

\geosection

\begin{frame}
  \frametitle{\subsecname}

  \begin{definition}
    A \textbf{crystal lattice} is a periodic arrangement of atoms
    repeating in three dimensions.
  \end{definition}
  \vfill
  Properties that follow from lattice symmetry:
  \vfill
  \begin{itemize}
    \item Cleavage planes along low-index directions
    \item Anisotropic optical and electrical behaviour
    \item Characteristic X-ray diffraction patterns
  \end{itemize}
\end{frame}

\subsection{The Seven Crystal Systems}

\begin{frame}
  \frametitle{\subsecname}
  \begin{columns}[T]
    \begin{column}{0.48\textwidth}
      \begin{itemize}
        \item Cubic
        \item Tetragonal
        \item Orthorhombic
        \item Hexagonal
      \end{itemize}
    \end{column}
    \begin{column}{0.48\textwidth}
      \begin{itemize}
        \item Trigonal
        \item Monoclinic
        \item Triclinic
      \end{itemize}
    \end{column}
  \end{columns}
\end{frame}

\section{A Coordination Example}

\geosection

\begin{frame}
  \frametitle{\subsecname}
  Coordination number in a simple cubic lattice:
  \vfill
  $$CN(\text{simple cubic}) = \textcolor{accent}{6}$$
  \vfill
  \begin{center}
    \begin{tikzpicture}[
      node/.style={circle, draw=accent!100, fill=bg1!40, very thick, minimum size=10mm},
      ]
      \node[node] (c) {A};
      \node[node] (n) [above=of c] {};
      \node[node] (s) [below=of c] {};
      \node[node] (e) [right=of c] {};
      \node[node] (w) [left=of c] {};

      \draw (c) -- (n);
      \draw (c) -- (s);
      \draw (c) -- (e);
      \draw (c) -- (w);
    \end{tikzpicture}
  \end{center}
\end{frame}

\section{Closing}

\begin{frame}
  \frametitle{Thank you}
  Questions welcome.\\[1em]
  dana.whitfield@example.org
\end{frame}

\end{document}
